\documentclass[11pt]{article}

\usepackage[T1]{fontenc}
\usepackage[utf8]{inputenc}
\usepackage[margin=1in]{geometry}
\usepackage{booktabs}
\usepackage{url}
\usepackage[hidelinks]{hyperref}

\title{GovernSpec: Zero-Intrusion Contract Compilation and Offline Acceptance Testing for Heterogeneous AI Agents}
\author{Ting Liu\\
SymbolicLight Research\\
Foshan, Guangdong, China\\
\texttt{research@symboliclight.com}}
\date{April 2026}

\begin{document}
\maketitle

\begin{center}
\small Project repository: \url{https://github.com/SymbolicLight-AGI/GovernSpec}
\end{center}

\begin{abstract}
AI-assisted software teams increasingly rely on heterogeneous agent interfaces:
repository instruction files, IDE rule files, structured-output schemas, and
machine-readable planning artifacts. The governance intent behind these artifacts is
often stable across tools, but the artifacts themselves are duplicated, edited, and
reviewed in incompatible formats. This fragmentation makes it difficult to keep
permissions, safety constraints, human confirmation rules, and output requirements
consistent without modifying every downstream agent runtime.

We present GovernSpec, a local-first tool that turns a single task-governance
contract into native artifacts for multiple AI-agent workflows. Developers author an
\texttt{govern.yaml} contract describing the task goal, permissions, constraints,
confirmation gates, output contract, and deterministic acceptance tests. GovernSpec
validates the contract, resolves reusable governance packs, normalizes it into an
intermediate intent representation, compiles target-specific artifacts, supports
reverse import from existing artifacts, and checks generated outputs with offline
assertions. The prototype currently implements 10 compile targets, 5 import source
types, and 10 deterministic assertion types. In a repository-local benchmark with
20 contracts, 52 output samples, and 20 handwritten artifacts, 94 of 120
representative contract-target pairs compiled successfully (78.33\%); compiled
round-trip import succeeded for 74 of 100 attempted artifacts (74.0\%);
handwritten artifact import succeeded for 20 of 20 artifacts; all 20 valid outputs
passed offline checks; and the targeted defect suite caught 32 of 32 labeled
defects. These results demonstrate
the feasibility of a zero-intrusion, artifact-level governance workflow for
settings in which teams cannot or do not want to modify agent runtimes.
\end{abstract}

\section{Introduction}

Large language models have become practical software engineering collaborators.
Code-specialized models and benchmarks show that language models can synthesize
programs, complete development tasks, and operate over real repositories (Chen et
al., 2021; Jimenez et al., 2024). Controlled and observational studies also suggest
that AI coding assistants can alter developer productivity and interaction patterns
(Peng et al., 2023; Ross et al., 2023). At the same time, the literature is cautious:
LLM-generated software can be wrong, insecure, or difficult to validate, and human
users may over-trust generated code (Fan et al., 2023; Perry et al., 2023). The
practical question for teams is therefore not only whether AI agents can help, but
how teams can make their intended constraints, permissions, and acceptance criteria
visible across the tools that agents already use.

Modern software teams rarely use a single AI-assistance surface. A repository may
contain durable project instructions, an IDE may consume project-scoped rule files,
an API integration may require a JSON schema, and an MCP-based workflow may exchange
machine-readable plans. These channels are useful because they meet developers where
they already work. They also create a governance maintenance problem: the same
requirements must be translated into several formats, each with different expressive
power and review affordances.

Consider a team that wants an agent to review a release plan. The policy is simple:
do not read private customer data, do not access the network, ask before destructive
file operations, produce a bounded report with required sections, and fail the
output if it contains unsupported claims. In practice, that policy may be copied
into markdown instructions, Cursor rule files, structured-output payloads, and MCP
plans. Each copy can drift. Some formats can express natural-language instructions
but not a schema. Some can express a JSON schema but not human confirmation gates. A
reviewer must understand every target format before judging whether the policy still
matches the original intent.

GovernSpec addresses this artifact-level problem. It is not an agent runtime, a
monitor, or a fail-closed policy enforcement layer. Instead, it provides a small
compiler and validator for task-governance contracts. A developer writes one
\texttt{govern.yaml} file. GovernSpec compiles that contract into the native
artifact channels used by downstream tools and then provides deterministic offline
checks for the final output.

This paper asks whether a single local contract can be compiled into heterogeneous
AI-agent artifacts, imported back into a structured representation, and used to
validate outputs without calling an external LLM service or modifying an agent
runtime.

The paper makes four concrete contributions:
\begin{itemize}
  \item A local-first contract workflow for making task goals, permissions,
  constraints, human gates, output requirements, and acceptance tests reviewable in
  one source.
  \item A target compiler that lowers the same contract into multiple agent-native
  artifacts, including instruction markdown, IDE rules, structured-output schemas,
  and MCP-style plans.
  \item A reverse-import workflow for migrating generated or handwritten artifacts
  back into draft contracts, exposing fidelity loss rather than hiding it.
  \item A reproducible artifact-level benchmark and offline assertion suite that
  evaluate compilation coverage, round-trip fidelity, and deterministic defect
  detection.
\end{itemize}

\section{Background and Positioning}

GovernSpec is motivated by a growing set of tool-specific artifact channels. Claude
Code documents project memory through \texttt{CLAUDE.md} files (Anthropic, n.d.).
Cursor project rules are stored as \texttt{.mdc} files with metadata and markdown
content (Cursor, n.d.). OpenAI Structured Outputs accept JSON Schema-based response
formats for schema-constrained model outputs (OpenAI, n.d.). The Model Context
Protocol defines a machine-readable protocol for connecting model applications with
tools and context sources (Model Context Protocol, n.d.).

These systems are not interchangeable. They expose different integration points,
different artifact formats, and different guarantees. GovernSpec therefore does not
attempt to define a new universal agent runtime. Its narrower contribution is a
compilation layer above existing artifact channels. The source contract captures the
governance intent once; target backends translate what each downstream channel can
represent; and capability notes make target limitations explicit.

This positioning also shapes the validation claim. Native instruction files and IDE
rules are advisory artifacts, not runtime monitors. GovernSpec therefore avoids
claiming that it can prevent unsafe actions at execution time. It demonstrates that
a team can author, propagate, migrate, and test governance intent locally, with
deterministic artifacts that are suitable for review and continuous integration.

\section{Related Work}

\subsection{LLMs for Software Engineering}

The starting point for GovernSpec is the rapid expansion of LLM-assisted software
engineering. Codex and HumanEval helped establish functional correctness as a
central evaluation concern for code generation (Chen et al., 2021). SWE-bench later
moved evaluation closer to real maintenance work by asking models to resolve GitHub
issues in existing repositories (Jimenez et al., 2024). Broader surveys of LLMs for
software engineering argue that LLMs are already relevant across coding, repair,
documentation, requirements, refactoring, and analytics, but also emphasize the need
for hybrid techniques that combine traditional software engineering checks with
model-generated work (Fan et al., 2023).

GovernSpec follows that hybrid direction. It does not attempt to improve the model
or to replace repository-level testing. Instead, it treats agent-facing artifacts as
software artifacts that should be parsed, compiled, reviewed, versioned, and tested.
This framing is deliberately conservative: the tool does not assume that a model
will obey every instruction, and it does not infer correctness from fluent output.

\subsection{Developer Interaction and AI Coding Assistants}

Studies of AI coding assistants show both practical benefit and operational risk.
Peng et al. (2023) report evidence that GitHub Copilot can improve developer
productivity in a controlled setting. Ross et al. (2023) study conversational
programming interactions with a large language model and highlight the importance of
interaction design around programming assistance. Perry et al. (2023), however,
show that users assisted by an AI coding system can produce less secure code while
also becoming more confident in their solutions.

These mixed findings motivate explicit task contracts. If a team relies on AI
assistance for code review, documentation, market analysis, or release planning, it
needs a maintainable way to state what the assistant is allowed to read, which
actions require human confirmation, and how outputs should be checked. GovernSpec
targets this coordination layer rather than the model layer.

\subsection{Prompt Engineering, Agents, and Tool Artifacts}

Prompt engineering research has shown that reusable prompt structures can improve
the controllability and repeatability of interactions with LLMs (Liu et al., 2023;
White et al., 2023). Agent research similarly studies systems in which language
models reason, call tools, and interact with environments (Xi et al., 2023; Yao et
al., 2023). These lines of work make instruction quality and tool context central to
agent behavior.

GovernSpec differs from a prompt catalog or an agent framework. Prompt patterns
describe reusable prompting strategies. Agent frameworks orchestrate model calls,
memory, tools, and actions. GovernSpec is a compilation and validation layer for the
artifacts around those systems. Its source contract can be compiled to prompt-like
markdown, IDE rules, structured-output schemas, or machine-readable plans, but the
tool intentionally avoids owning the runtime loop.

\subsection{Structured Data and Policy-as-Code}

Structured output channels draw on a longer tradition of machine-readable data
formats and schemas. JSON is standardized as a lightweight data interchange format
(Bray, 2017), and JSON Schema defines vocabularies for describing and validating
JSON documents (Wright et al., 2022). GovernSpec uses this ecosystem when compiling
structured-output targets and when checking JSON outputs offline.

GovernSpec is also adjacent to policy-as-code. Policy-as-code systems encode
governance rules in machine-readable artifacts so they can be reviewed, tested, and
applied in automated workflows. Recent empirical work studies how policy-as-code is
adopted and maintained in open-source repositories (Foalem et al., 2026). GovernSpec
borrows the same engineering instinct but applies it to AI-agent task artifacts. The
important distinction is that GovernSpec is not a runtime authorization engine. It
is a local contract compiler and acceptance checker for artifacts that may feed
different agent tools.

\subsection{Adjacent Intent Validators}

Lightweight CI validators address a narrower part of the same artifact-management
space. For example, JanneL's \texttt{validate-intentspec-action} validates
Markdown front matter against a small schema inside GitHub Actions. GovernSpec uses
a different source format and scope: YAML contracts, reusable governance packs, a
normalized intermediate representation, target-specific compilation, reverse import,
and offline output assertions. This makes the projects complementary rather than
interchangeable; the action checks whether one Markdown artifact has required
front-matter fields, while GovernSpec treats the source contract as an authoring,
compilation, migration, and testing object.

\section{Design Goals}

GovernSpec is organized around five design goals.

\paragraph{Local-first reproducibility.}
All core workflows should run without remote model calls, API keys, or network
access. This makes the tool suitable for repositories that handle private data and
for continuous integration environments where external calls are restricted.

\paragraph{Zero-intrusion adoption.}
Teams should not need to modify an agent runtime before they can benefit from a
single source of governance intent. GovernSpec therefore emits artifacts that
existing tools can already consume.

\paragraph{Reviewable governance intent.}
The source contract should make task goals, allowed and forbidden capabilities,
human gates, output obligations, and acceptance tests visible in a compact file that
can be reviewed like code.

\paragraph{Explicit target loss.}
Not every target can represent every field. A markdown instruction file can express
many natural-language constraints but cannot enforce a JSON schema. A
structured-output schema can constrain the output shape but cannot represent all
permissions or human gates. GovernSpec records these losses as capability notes.

\paragraph{Deterministic offline checks.}
The test runner should evaluate observable output properties deterministically. It
should not claim to judge semantic truth, model intent, or runtime obedience.

\section{Tool Design}

GovernSpec has five internal stages:
\begin{verbatim}
govern.yaml
  -> parse and validate
  -> resolve imported governance packs
  -> normalize to IIR
  -> compile to target-specific artifacts
  -> test generated outputs offline
\end{verbatim}

The source contract is a YAML document containing metadata, task context, inputs,
permissions, constraints, evidence expectations, output requirements, human
confirmation gates, and tests. Reusable policy fragments are represented as
\texttt{GovernPack} files. Import resolution applies conservative merge behavior,
including deny-wins permissions and preservation of imported acceptance assertions.

The intermediate intent representation (IIR) separates authoring concerns from
target formatting. It stores the normalized task goal, resolved permissions, merged
constraints, output contract, human gates, test contracts, risk signals, and target
capability notes. Backends then lower the IIR into concrete target families. For
instruction-style targets, GovernSpec emits markdown. For Cursor rules, it emits a
rule bundle. For structured-output targets, it emits JSON schema payloads. For MCP
planning, it emits a machine-readable plan with risk and constraint-loss fields.

The offline acceptance runner evaluates final artifacts after an agent has produced
an output. It supports required sections, required and forbidden substrings, regular
expression checks, word and character limits, JSON schema validation, JSON path
existence, and JSON array size constraints. These checks deliberately focus on
deterministic compliance properties. They do not verify semantic truth, factual
correctness, or whether a model internally followed the instructions.

\section{Demonstration Scenario}

The demonstration uses the repository-local benchmark in
\path{benchmark/reproducible_artifact/}. It contains 20 contracts, 52 output
samples, 20 handwritten artifacts, explicit labels, experiment scripts, and
generated results.
The benchmark is designed to run without external model APIs or network-dependent
agent services.

The live demonstration has four segments. First, the presenter authors or edits a
small YAML contract with a task goal, permissions, human gates, output
requirements, and tests. Second, GovernSpec compiles the same contract into
instruction markdown, Cursor rules, OpenAI and Gemini structured-output payloads,
and an MCP plan. Third, the presenter reverse-imports generated and handwritten
artifacts into draft contracts and shows where fidelity is preserved or lost.
Fourth, the presenter runs the offline acceptance runner against valid and defective
outputs, showing deterministic failure messages for missing sections, forbidden
text, malformed JSON, and schema mismatch.

All experiments are reproducible with a single command from the repository root:
\begin{quote}
\small\texttt{python} \path{benchmark/reproducible_artifact/scripts/run_all.py}
\end{quote}
The command generates JSON result files and paper-friendly Markdown summaries for
the compile matrix, round-trip fidelity, assertion evaluation, aggregate numbers,
and rendered tables.

\section{Evaluation Design}

The validation is intentionally artifact-level. It evaluates whether the current
implementation can compile representative contracts, recover structured drafts from
artifacts, and detect labeled output defects. It does not evaluate live agent
behavior.

\subsection{Benchmark Assets}

The benchmark contains 20 contracts. Five are adapted from existing project examples
and 15 are paper-only contracts designed to exercise additional output and risk
patterns. The contracts cover markdown and JSON outputs, English and Chinese
content, imported governance packs, human gates, and representative risks involving
network access, confidential inputs, and filesystem or tool use. The benchmark also
contains 20 valid outputs, 32 targeted defect outputs, and 20 handwritten artifacts
that resemble existing repository instruction files or structured-output payloads.

The defect outputs are labeled by targeted assertion type. Each isolated defect
sample is intended to violate one supported assertion type while keeping the rest of
the output close to a valid example. This design is small, but it makes the failure
mode inspectable and reproducible.

\subsection{Metrics}

The compile matrix reports whether each contract-target pair compiles, the artifact
kind, error type, error message, and the number of target capability notes. The
round-trip experiment compiles contracts into importable targets, imports the
artifacts back, and compares the recovered drafts on core fields:
\texttt{goal}, \texttt{permissions}, \texttt{constraints},
\texttt{human\_gates}, \texttt{output}, and \texttt{tests}. Metadata and source-path
fields are ignored. JSON schemas are normalized with stable key ordering before
comparison.

The assertion experiment reports whether each valid or defective sample has the
expected acceptance outcome. For targeted defects, a caught defect means the
offline runner returns failure for a sample labeled as expected to fail. The summary
also reports per-assertion coverage.

\section{Results}

\subsection{Multi-Target Compilation}

The compile matrix evaluates 20 contracts against 6 representative targets, producing
120 contract-target pairs. Table~\ref{tab:compile} reports target-level coverage.
Instruction-style targets and \texttt{mcp-plan} compiled for all contracts.
Structured-output targets compiled for JSON-output contracts and failed for
markdown-output contracts. These failures are expected target limitations because
schema-constrained output targets require JSON output contracts. Overall, 94 of 120
contract-target pairs compiled successfully (78.33\%).

\begin{table}[htbp]
\caption{Target coverage in the compile matrix.}
\label{tab:compile}
\centering
\begin{tabular}{lrr}
\toprule
Target & Succeeded & Total \\
\midrule
\texttt{agents-md} & 20 & 20 \\
\texttt{claude-md} & 20 & 20 \\
\texttt{cursor-rules} & 20 & 20 \\
\texttt{openai-structured} & 7 & 20 \\
\texttt{gemini-structured} & 7 & 20 \\
\texttt{mcp-plan} & 20 & 20 \\
\bottomrule
\end{tabular}
\end{table}

The important observation is not that every target succeeds. Rather, GovernSpec
distinguishes between successful compilation and target-specific unsuitability. A
markdown report contract can be represented as instructions, rules, and planning
metadata, but it cannot be faithfully lowered into a schema-constrained JSON
response format without changing the intended output kind. Recording this as a
compile limitation is preferable to silently producing a misleading schema.

\subsection{Reverse Import Fidelity}

The round-trip experiment compiles contracts into the 5 importable source types and
imports them back into draft contracts. The comparison focuses on core fields:
\texttt{goal}, \texttt{permissions}, \texttt{constraints}, \texttt{human\_gates},
\texttt{output}, and \texttt{tests}. Non-core metadata and source-path fields are
ignored. JSON schemas are normalized with stable key ordering before comparison.

\begin{table}[htbp]
\caption{Round-trip import and core-field fidelity.}
\label{tab:roundtrip}
\centering
\begin{tabular}{lrr}
\toprule
Source type & Imported & Fidelity \\
\midrule
\texttt{agents-md} & 20 & 77.5\% \\
\texttt{claude-md} & 20 & 77.5\% \\
\texttt{cursor-rules} & 20 & 77.5\% \\
\texttt{gemini-structured} & 7 & 47.62\% \\
\texttt{openai-structured} & 7 & 47.62\% \\
\bottomrule
\end{tabular}
\end{table}

Compiled round-trip import succeeded for 74 of 100 attempted artifacts (74.0\%).
Successful compiled round trips had an average core-field fidelity of 71.85\%. The
benchmark also imports 20 handwritten artifacts and matches them against partial gold
labels; all 20 imported successfully. These results show that generated instruction
artifacts preserve more recoverable governance information than structured-output
payloads, while structured-output payloads remain useful for output schema recovery.

\subsection{Offline Assertion Effectiveness}

The assertion experiment evaluates 20 valid outputs and 32 targeted defect outputs.
The defect suite contains at least two samples for each supported assertion type.
All 20 valid outputs passed offline acceptance tests (100.0\%), and the targeted
defect suite caught 32 of 32 labeled defects (100.0\%).

\begin{table}[htbp]
\caption{Targeted defect detection by assertion type.}
\label{tab:assertions}
\centering
\begin{tabular}{lrrr}
\toprule
Assertion & Samples & Caught & Rate \\
\midrule
\texttt{contains} & 6 & 6 & 100.0\% \\
\texttt{json\_array\_min\_items} & 2 & 2 & 100.0\% \\
\texttt{json\_path\_exists} & 2 & 2 & 100.0\% \\
\texttt{json\_schema} & 2 & 2 & 100.0\% \\
\texttt{max\_chars} & 2 & 2 & 100.0\% \\
\texttt{max\_words} & 3 & 3 & 100.0\% \\
\texttt{no\_regex} & 4 & 4 & 100.0\% \\
\texttt{not\_contains} & 4 & 4 & 100.0\% \\
\texttt{regex} & 3 & 3 & 100.0\% \\
\texttt{required\_sections} & 4 & 4 & 100.0\% \\
\bottomrule
\end{tabular}
\end{table}

The observed failures cover missing markdown sections, forbidden content, absent
required content, regex mismatch, forbidden regex matches, length violations, JSON
schema mismatch, missing JSON paths, and undersized JSON arrays. The suite is not a
semantic evaluator. It is a deterministic guardrail for properties that teams can
state before running an agent and check after receiving an output.

\section{Case Studies}

\subsection{Multi-Target Governance Propagation}

One benchmark contract describes a code-review task with restricted permissions,
required evidence, bounded output sections, and acceptance checks over markdown. The
same source contract compiles into repository instructions, Claude-style memory,
Cursor rules, and an MCP plan. The instruction targets preserve most governance
content as reviewable text. The MCP plan additionally exposes risk and
constraint-loss fields, making it useful for inspecting what a downstream planning
workflow can and cannot represent.

This case study illustrates the main authoring benefit. A policy reviewer can
inspect one source contract and generated target artifacts, rather than manually
checking several independently edited files. The generated artifacts do not make the
agent safe by themselves, but they reduce drift in the instructions supplied to
different tools.

\subsection{Reverse Import for Migration}

Another case study starts from existing handwritten artifacts. These artifacts
include instruction markdown, Claude-style memory, Cursor rules, and structured
payloads. GovernSpec imports each artifact into a draft contract and compares the
result with partial gold labels. The handwritten import experiment succeeds for all
6 artifacts.

The value of this workflow is migration rather than perfect reconstruction.
Existing teams often have useful instructions before they have a formal contract
format. Reverse import gives them a draft starting point and highlights missing
fields, such as explicit permissions or output tests, that were implicit in the
original artifact.

\subsection{Defect Detection Before Publication}

The assertion suite provides a third case study. For each supported assertion type,
the benchmark includes multiple targeted defective outputs. The runner catches all 32
labeled defects. Examples include a markdown report missing a required section, a
JSON output missing a required path, a response that exceeds a word limit, and an
output containing forbidden text.

This case study is intentionally modest. It does not show that the output is true or
that the agent behaved correctly. It shows that teams can express low-level
acceptance criteria in the same contract used to generate agent artifacts, and then
apply those criteria deterministically in local automation.

\section{Discussion}

\subsection{Why Compilation Instead of Runtime Enforcement}

Runtime enforcement is valuable when a system controls the execution boundary. Many
teams, however, adopt agent tools whose internal runtime loops are not controlled by
the repository owner. In those settings, a zero-intrusion compilation layer can be a
useful intermediate step. It does not replace sandboxing, access control, or policy
engines, but it improves the maintainability and reviewability of the artifacts that
developers already ship to agents.

\subsection{Why Reverse Import Matters}

Governance artifacts rarely begin life as formal schemas. They often start as
repository notes, project conventions, IDE rules, or copied prompt fragments.
Reverse import acknowledges this reality. By recovering a draft contract from
existing artifacts, GovernSpec supports incremental adoption. The fidelity score
also keeps migration honest: if an artifact does not contain explicit tests or
human gates, the imported contract should not pretend that it recovered them.

\subsection{Interpreting the Benchmark}

The benchmark should be read as a reproducibility and feasibility artifact, not as a
claim of broad empirical generality. Its strongest result is that the pipeline can
run locally end to end and expose target limitations in a structured way. Its
weakest area is scale. Future work should add more repositories, more handwritten
artifacts, independent labels, and live-agent studies that measure whether compiled
artifacts affect model behavior in practice.

\section{Threats to Validity and Limitations}

\paragraph{External validity.}
The benchmark is small and curated. It is appropriate for demonstrating feasibility,
but it is not evidence of broad empirical generality across all agent tools,
repositories, or policy styles. The contracts intentionally cover markdown and JSON
outputs, English and Chinese text, imports, human gates, filesystem and network
risks, and all supported assertion types, but the suite remains a seed benchmark.

\paragraph{Construct validity.}
The validation isolates artifact behavior from model behavior. This design makes the
experiments deterministic and reproducible, but it also means the results do not
measure whether a live agent follows a compiled instruction file in practice.
GovernSpec should therefore be viewed as an authoring, propagation, migration, and
post-hoc validation layer rather than as runtime security enforcement.

\paragraph{Internal validity.}
The current labels were prepared for a demonstration benchmark rather than through a
multi-annotator study. This is acceptable for an initial artifact evaluation, but
future studies should include independent labeling and agreement statistics.

\paragraph{Reverse-import limitations.}
Reverse import is heuristic for natural-language artifacts. It can recover goals,
constraints, permissions, and output hints from generated or handwritten artifacts,
but it cannot guarantee semantic equivalence with the original author intent. The
tool exposes this limitation through field-level fidelity rather than hiding it
behind a binary success metric.

\section{Availability}

GovernSpec is implemented as a Python 3.11+ local CLI and library. The benchmark and
generated results are included in the public repository at
\url{https://github.com/SymbolicLight-AGI/GovernSpec}, under
\path{benchmark/reproducible_artifact/}. The prototype does not call real LLM APIs,
does not require API keys, and does not require network access for the reported
experiments. The local reproduction command is
\texttt{python} \path{benchmark/reproducible_artifact/scripts/run_all.py}.

\section*{References}

Anthropic. (n.d.). \textit{How Claude remembers your project}. Claude Code Docs.
Retrieved April 24, 2026, from
\url{https://code.claude.com/docs/en/memory}

Bray, T. (2017). \textit{The JavaScript Object Notation (JSON) data interchange
format} (RFC 8259). Internet Engineering Task Force.
\url{https://www.rfc-editor.org/rfc/rfc8259}

Chen, M., Tworek, J., Jun, H., Yuan, Q., Pinto, H. P. de O., Kaplan, J.,
Edwards, H., Burda, Y., Joseph, N., Brockman, G., Ray, A., Puri, R.,
Krueger, G., Petrov, M., Khlaaf, H., Sastry, G., Mishkin, P., Chan, B.,
Gray, S., ... Zaremba, W. (2021). \textit{Evaluating large language models
trained on code}. arXiv. \url{https://arxiv.org/abs/2107.03374}

Cursor. (n.d.). \textit{Rules}. Cursor Documentation. Retrieved April 24, 2026,
from \url{https://docs.cursor.com/context/rules}

Fan, A., Gokkaya, B., Harman, M., Lyubarskiy, M., Sengupta, S., Yoo, S., \&
Zhang, J. M. (2023). \textit{Large language models for software engineering:
Survey and open problems}. arXiv. \url{https://arxiv.org/abs/2310.03533}

Foalem, P. L., Khomh, F., Da Silva, L., \& Merlo, E. (2026).
\textit{An empirical study of policy-as-code adoption in open-source software
projects}. arXiv. \url{https://arxiv.org/abs/2601.05555}

JanneL. (2026). \textit{Validate IntentSpec Action}. GitHub. Retrieved May 11,
2026, from \url{https://github.com/JanneL/validate-intentspec-action}

Jimenez, C. E., Yang, J., Wettig, A., Yao, S., Pei, K., Press, O., \&
Narasimhan, K. R. (2024). \textit{SWE-bench: Can language models resolve
real-world GitHub issues?} International Conference on Learning Representations.
\url{https://openreview.net/forum?id=VTF8yNQM66}

Liu, P., Yuan, W., Fu, J., Jiang, Z., Hayashi, H., \& Neubig, G. (2023).
Pre-train, prompt, and predict: A systematic survey of prompting methods in
natural language processing. \textit{ACM Computing Surveys, 55}(9), Article
195. \url{https://doi.org/10.1145/3560815}

Model Context Protocol. (n.d.). \textit{Specification}. Model Context Protocol
Documentation. Retrieved April 24, 2026, from
\url{https://modelcontextprotocol.io/specification/draft}

OpenAI. (n.d.). \textit{Structured model outputs}. OpenAI API Documentation.
Retrieved April 24, 2026, from
\url{https://developers.openai.com/api/docs/guides/structured-outputs}

Peng, S., Kalliamvakou, E., Cihon, P., \& Demirer, M. (2023). \textit{The
impact of AI on developer productivity: Evidence from GitHub Copilot}. arXiv.
\url{https://arxiv.org/abs/2302.06590}

Perry, N., Srivastava, M., Kumar, D., \& Boneh, D. (2023). Do users write more
insecure code with AI assistants? In \textit{Proceedings of the 2023 ACM SIGSAC
Conference on Computer and Communications Security} (pp. 2785--2799). ACM.
\url{https://doi.org/10.1145/3576915.3623157}

Ross, S. I., Martinez, F., Houde, S., Muller, M., \& Weisz, J. D. (2023). The
programmer's assistant: Conversational interaction with a large language model
for software development. In \textit{Proceedings of the 28th International
Conference on Intelligent User Interfaces} (pp. 491--514). ACM.
\url{https://doi.org/10.1145/3581641.3584037}

White, J., Fu, Q., Hays, S., Sandborn, M., Olea, C., Gilbert, H., Elnashar, A.,
Spencer-Smith, J., \& Schmidt, D. C. (2023). \textit{A prompt pattern catalog
to enhance prompt engineering with ChatGPT}. arXiv.
\url{https://arxiv.org/abs/2302.11382}

Wright, A., Andrews, H., Hutton, B., \& Dennis, G. (2022). \textit{JSON Schema
Draft 2020-12}. JSON Schema. \url{https://json-schema.org/draft/2020-12}

Xi, Z., Chen, W., Guo, X., He, W., Ding, Y., Hong, B., Zhang, M., Wang, J.,
Jin, S., Zhou, E., Zheng, R., Fan, X., Wang, X., Xiong, L., Zhou, Y., Wang, W.,
Jiang, C., Zou, Y., Liu, X., ... Gui, T. (2023). \textit{The rise and potential
of large language model based agents: A survey}. arXiv.
\url{https://arxiv.org/abs/2309.07864}

Yao, S., Zhao, J., Yu, D., Du, N., Shafran, I., Narasimhan, K. R., \& Cao, Y.
(2023). ReAct: Synergizing reasoning and acting in language models. In
\textit{International Conference on Learning Representations}.
\url{https://openreview.net/forum?id=WE_vluYUL-X}

\end{document}
